\documentclass[11pt,a4paper]{article}
\usepackage[margin=1in]{geometry}\usepackage[T1]{fontenc}\usepackage[utf8]{inputenc}
\usepackage{enumitem,xcolor,hyperref,booktabs}
\newcommand{\nRawRows}{25,396}
\newcommand{\nPlayerSeasons}{20,547}
\newcommand{\nAllPlayers}{3,817}
\newcommand{\nDebutWindow}{2,381}
\newcommand{\nGamesFortytwo}{1,833}
\newcommand{\nPlayers}{1,676}
\newcommand{\nObs}{13,193}
\newcommand{\nActive}{17}
\newcommand{\minActiveSeason}{15}
\newcommand{\sdPerLowMp}{14.9}
\newcommand{\sdPerHighMp}{4.1}
\newcommand{\minSpan}{2}
\newcommand{\maxSpan}{23}
\newcommand{\nGapPlayers}{458}
\newcommand{\pctGapPlayers}{27.3}
\newcommand{\nGapSeasons}{798}
\newcommand{\nSpanSeasons}{13,991}
\newcommand{\pctGapSeasons}{5.7}
\newcommand{\muPer}{13.67}
\newcommand{\sdPer}{4.32}
\newcommand{\skewPer}{0.39}
\newcommand{\nDev}{1,341}
\newcommand{\nHeld}{335}
\newcommand{\nDrafted}{1,336}
\newcommand{\medianSeasons}{7}
\newcommand{\meanSeasons}{7.9}
\newcommand{\nGrid}{720}
\newcommand{\pcaVar}{96}
\newcommand{\rmseHeld}{0.74}
\newcommand{\rmseHeldSd}{0.06}
\newcommand{\rmseTrain}{0.43}
\newcommand{\rmseLinear}{2.24}
\newcommand{\rmsePlayerMean}{2.76}
\newcommand{\rmseGlobal}{4.35}
\newcommand{\rmseHeldEight}{1.23}
\newcommand{\rmseHeldSixteen}{0.95}
\newcommand{\rmseHeldThirtytwo}{0.80}
\newcommand{\rmseHeldOneTwoEight}{0.74}
\newcommand{\bestEpoch}{361}
\newcommand{\nParams}{179,073}
\newcommand{\rmseOld}{2.23}
\newcommand{\selDbcvKdev}{7}
\newcommand{\selDbcvKho}{7}
\newcommand{\selDbcvKall}{13}
\newcommand{\selDbcvNoiseho}{1}
\newcommand{\selSilKdev}{24}
\newcommand{\selSilKho}{23}
\newcommand{\selSilKall}{14}
\newcommand{\selSilNoiseho}{56}
\newcommand{\selCompKdev}{14}
\newcommand{\selCompKho}{14}
\newcommand{\selCompKall}{7}
\newcommand{\selCompNoiseho}{29}
\newcommand{\gridKmin}{2}
\newcommand{\gridKmax}{33}
\newcommand{\ariSingleMin}{0.32}
\newcommand{\ariSingleMax}{0.66}
\newcommand{\ariAeSeedsMin}{0.09}
\newcommand{\ariAeSeedsMax}{0.38}
\newcommand{\kSingleMin}{2}
\newcommand{\kSingleMax}{27}
\newcommand{\bigShareEomMax}{63}
\newcommand{\bigShareEomMin}{33}
\newcommand{\noiseLeafMean}{34}
\newcommand{\noiseEomMean}{14}
\newcommand{\configList}{C1: 10,10,0.0,25,15; C2: 5,15,0.0,30,10; C3: 5,20,0.0,22,7; C4: 2,15,0.0,40,5}
\newcommand{\aeAllK}{3}
\newcommand{\aeAllSil}{0.48}
\newcommand{\aeAllPac}{0.57}
\newcommand{\aeAllStable}{80}
\newcommand{\aeAllSilEight}{0.44}
\newcommand{\aeAllPacEight}{0.32}
\newcommand{\aeAllStableEight}{67}
\newcommand{\aeAllRsqEight}{0.50}
\newcommand{\aeAllRsq}{0.37}
\newcommand{\aeSeedZeroK}{3}
\newcommand{\aeSeedZeroSil}{0.90}
\newcommand{\aeSeedZeroPac}{0.17}
\newcommand{\aeSeedZeroStable}{100}
\newcommand{\aeSeedZeroSilEight}{0.78}
\newcommand{\aeSeedZeroPacEight}{0.13}
\newcommand{\aeSeedZeroStableEight}{95}
\newcommand{\aeSeedZeroRsqEight}{0.50}
\newcommand{\aeSeedZeroRsq}{0.41}
\newcommand{\aeNostdK}{3}
\newcommand{\aeNostdSil}{0.56}
\newcommand{\aeNostdPac}{0.50}
\newcommand{\aeNostdStable}{94}
\newcommand{\aeNostdSilEight}{0.45}
\newcommand{\aeNostdPacEight}{0.31}
\newcommand{\aeNostdStableEight}{70}
\newcommand{\aeNostdRsqEight}{0.57}
\newcommand{\aeNostdRsq}{0.46}
\newcommand{\aeUmapK}{3}
\newcommand{\aeUmapSil}{0.71}
\newcommand{\aeUmapPac}{0.38}
\newcommand{\aeUmapStable}{95}
\newcommand{\aeUmapSilEight}{0.38}
\newcommand{\aeUmapPacEight}{0.33}
\newcommand{\aeUmapStableEight}{49}
\newcommand{\aeUmapRsqEight}{0.29}
\newcommand{\aeUmapRsq}{0.24}
\newcommand{\aeHdbK}{5}
\newcommand{\aeHdbSil}{0.43}
\newcommand{\aeHdbPac}{0.76}
\newcommand{\aeHdbStable}{66}
\newcommand{\aeHdbSilEight}{0.36}
\newcommand{\aeHdbPacEight}{0.76}
\newcommand{\aeHdbStableEight}{66}
\newcommand{\aeHdbRsqEight}{0.18}
\newcommand{\aeHdbRsq}{0.18}
\newcommand{\summK}{4}
\newcommand{\summSil}{0.93}
\newcommand{\summPac}{0.08}
\newcommand{\summStable}{99}
\newcommand{\summSilEight}{0.90}
\newcommand{\summPacEight}{0.07}
\newcommand{\summStableEight}{99}
\newcommand{\summRsqEight}{0.60}
\newcommand{\summRsq}{0.47}
\newcommand{\summWardK}{3}
\newcommand{\summWardSil}{0.49}
\newcommand{\summWardPac}{0.62}
\newcommand{\summWardStable}{96}
\newcommand{\summWardSilEight}{0.47}
\newcommand{\summWardPacEight}{0.33}
\newcommand{\summWardStableEight}{81}
\newcommand{\summWardRsqEight}{0.55}
\newcommand{\summWardRsq}{0.42}
\newcommand{\rawK}{3}
\newcommand{\rawSil}{0.93}
\newcommand{\rawPac}{0.10}
\newcommand{\rawStable}{99}
\newcommand{\rawSilEight}{0.85}
\newcommand{\rawPacEight}{0.10}
\newcommand{\rawStableEight}{97}
\newcommand{\rawRsqEight}{0.72}
\newcommand{\rawRsq}{0.57}
\newcommand{\pcaK}{3}
\newcommand{\pcaSil}{0.74}
\newcommand{\pcaPac}{0.31}
\newcommand{\pcaStable}{94}
\newcommand{\pcaSilEight}{0.50}
\newcommand{\pcaPacEight}{0.24}
\newcommand{\pcaStableEight}{80}
\newcommand{\pcaRsqEight}{0.34}
\newcommand{\pcaRsq}{0.14}
\newcommand{\dtwK}{3}
\newcommand{\dtwSil}{0.76}
\newcommand{\dtwPac}{0.34}
\newcommand{\dtwStable}{97}
\newcommand{\dtwSilEight}{0.36}
\newcommand{\dtwPacEight}{0.35}
\newcommand{\dtwStableEight}{70}
\newcommand{\dtwRsqEight}{0.66}
\newcommand{\dtwRsq}{0.53}
\newcommand{\aeEightK}{3}
\newcommand{\aeEightSil}{0.50}
\newcommand{\aeEightPac}{0.55}
\newcommand{\aeEightStable}{90}
\newcommand{\aeEightSilEight}{0.34}
\newcommand{\aeEightPacEight}{0.35}
\newcommand{\aeEightStableEight}{46}
\newcommand{\aeEightRsqEight}{0.48}
\newcommand{\aeEightRsq}{0.40}
\newcommand{\aeSixteenK}{4}
\newcommand{\aeSixteenSil}{0.57}
\newcommand{\aeSixteenPac}{0.41}
\newcommand{\aeSixteenStable}{86}
\newcommand{\aeSixteenSilEight}{0.35}
\newcommand{\aeSixteenPacEight}{0.32}
\newcommand{\aeSixteenStableEight}{53}
\newcommand{\aeSixteenRsqEight}{0.48}
\newcommand{\aeSixteenRsq}{0.38}
\newcommand{\aeThirtytwoK}{3}
\newcommand{\aeThirtytwoSil}{0.62}
\newcommand{\aeThirtytwoPac}{0.46}
\newcommand{\aeThirtytwoStable}{91}
\newcommand{\aeThirtytwoSilEight}{0.37}
\newcommand{\aeThirtytwoPacEight}{0.33}
\newcommand{\aeThirtytwoStableEight}{50}
\newcommand{\aeThirtytwoRsqEight}{0.49}
\newcommand{\aeThirtytwoRsq}{0.39}
\newcommand{\aeOneTwoEightK}{3}
\newcommand{\aeOneTwoEightSil}{0.70}
\newcommand{\aeOneTwoEightPac}{0.39}
\newcommand{\aeOneTwoEightStable}{92}
\newcommand{\aeOneTwoEightSilEight}{0.49}
\newcommand{\aeOneTwoEightPacEight}{0.29}
\newcommand{\aeOneTwoEightStableEight}{71}
\newcommand{\aeOneTwoEightRsqEight}{0.46}
\newcommand{\aeOneTwoEightRsq}{0.39}
\newcommand{\summSilTwo}{0.98}
\newcommand{\summSilFour}{0.93}
\newcommand{\summSilTwelve}{0.59}
\newcommand{\aeAllSilThree}{0.48}
\newcommand{\aeAllSilTwelve}{0.34}
\newcommand{\aeAllPacThree}{0.57}
\newcommand{\nMembersAe}{40}
\newcommand{\nMembersSumm}{8}
\newcommand{\mainK}{4}
\newcommand{\fourNA}{369}
\newcommand{\fourPerA}{16.5}
\newcommand{\fourSeasA}{12}
\newcommand{\fourASA}{52}
\newcommand{\fourStabA}{0.98}
\newcommand{\fourNB}{560}
\newcommand{\fourPerB}{12.5}
\newcommand{\fourSeasB}{11}
\newcommand{\fourASB}{8}
\newcommand{\fourStabB}{0.98}
\newcommand{\fourNC}{444}
\newcommand{\fourPerC}{10.6}
\newcommand{\fourSeasC}{3}
\newcommand{\fourASC}{1}
\newcommand{\fourStabC}{0.96}
\newcommand{\fourND}{303}
\newcommand{\fourPerD}{10.6}
\newcommand{\fourSeasD}{3}
\newcommand{\fourASD}{0}
\newcommand{\fourStabD}{0.98}
\newcommand{\eightNA}{138}
\newcommand{\eightPerA}{18.5}
\newcommand{\eightSeasA}{14}
\newcommand{\eightASA}{86}
\newcommand{\eightStabA}{0.98}
\newcommand{\eightNB}{166}
\newcommand{\eightPerB}{15.1}
\newcommand{\eightSeasB}{6}
\newcommand{\eightASB}{19}
\newcommand{\eightStabB}{0.95}
\newcommand{\eightNC}{283}
\newcommand{\eightPerC}{14.2}
\newcommand{\eightSeasC}{13}
\newcommand{\eightASC}{28}
\newcommand{\eightStabC}{0.96}
\newcommand{\eightND}{141}
\newcommand{\eightPerD}{13.8}
\newcommand{\eightSeasD}{7}
\newcommand{\eightASD}{6}
\newcommand{\eightStabD}{0.96}
\newcommand{\eightNE}{276}
\newcommand{\eightPerE}{11.2}
\newcommand{\eightSeasE}{4}
\newcommand{\eightASE}{0}
\newcommand{\eightStabE}{0.95}
\newcommand{\eightNF}{339}
\newcommand{\eightPerF}{11.2}
\newcommand{\eightSeasF}{10}
\newcommand{\eightASF}{1}
\newcommand{\eightStabF}{0.95}
\newcommand{\eightNG}{206}
\newcommand{\eightPerG}{9.5}
\newcommand{\eightSeasG}{3}
\newcommand{\eightASG}{0}
\newcommand{\eightStabG}{0.98}
\newcommand{\eightNH}{127}
\newcommand{\eightPerH}{9.5}
\newcommand{\eightSeasH}{3}
\newcommand{\eightASH}{0}
\newcommand{\eightStabH}{0.90}
\newcommand{\ariFourEight}{0.37}
\newcommand{\nestingPurity}{79}
\newcommand{\typeJordan}{T1}
\newcommand{\stabJordan}{0.98}
\newcommand{\typeNash}{T1}
\newcommand{\stabNash}{0.99}
\newcommand{\typeSampson}{T2}
\newcommand{\stabSampson}{0.95}
\newcommand{\typeMilicic}{T6}
\newcommand{\stabMilicic}{0.95}
\newcommand{\typeMorrison}{T5}
\newcommand{\stabMorrison}{0.96}
\newcommand{\typeBryant}{T1}
\newcommand{\stabBryant}{0.98}
\newcommand{\typeRose}{T1}
\newcommand{\stabRose}{0.99}
\newcommand{\typeIsiah}{T1}
\newcommand{\stabIsiah}{0.98}
\newcommand{\typeLeBron}{T1}
\newcommand{\stabLeBron}{0.98}
\newcommand{\typeCurry}{T1}
\newcommand{\stabCurry}{0.98}
\newcommand{\typeBird}{T1}
\newcommand{\stabBird}{0.98}
\newcommand{\typeGinobili}{T1}
\newcommand{\stabGinobili}{0.98}
\newcommand{\typeWallace}{T3}
\newcommand{\stabWallace}{0.97}
\newcommand{\VPos}{0.08}
\newcommand{\pPos}{0.086}
\newcommand{\amiPos}{0.00}
\newcommand{\VDraft}{0.27}
\newcommand{\pDraft}{0.000}
\newcommand{\amiDraft}{0.06}
\newcommand{\VAllStar}{0.68}
\newcommand{\pAllStar}{0.000}
\newcommand{\amiAllStar}{0.16}
\newcommand{\VAllNba}{0.68}
\newcommand{\pAllNba}{0.000}
\newcommand{\amiAllNba}{0.12}
\newcommand{\epsWs}{0.75}
\newcommand{\epsBpm}{0.54}
\newcommand{\epsGames}{0.74}
\newcommand{\epsPer}{0.75}
\newcommand{\epsLen}{0.76}
\newcommand{\ariAAeSeedOnly}{0.58}
\newcommand{\ariAEightAeSeedOnly}{0.57}
\newcommand{\kAAeSeedOnly}{3}
\newcommand{\ariAAeNostd}{0.54}
\newcommand{\ariAEightAeNostd}{0.52}
\newcommand{\kAAeNostd}{3}
\newcommand{\ariAAeUmap}{0.31}
\newcommand{\ariAEightAeUmap}{0.29}
\newcommand{\kAAeUmap}{3}
\newcommand{\ariAAeUmapHdbscan}{0.24}
\newcommand{\ariAEightAeUmapHdbscan}{0.18}
\newcommand{\kAAeUmapHdbscan}{5}
\newcommand{\ariAAe}{0.62}
\newcommand{\ariAEightAe}{0.53}
\newcommand{\kAAe}{3}
\newcommand{\ariAAeOriginal}{0.20}
\newcommand{\ariAEightAeOriginal}{0.18}
\newcommand{\kAAeOriginal}{3}
\newcommand{\ariAAeMp}{0.35}
\newcommand{\ariAEightAeMp}{0.29}
\newcommand{\kAAeMp}{3}
\newcommand{\ariAAeMinSeasons}{0.35}
\newcommand{\ariAEightAeMinSeasons}{0.27}
\newcommand{\kAAeMinSeasons}{3}
\newcommand{\ariAAeNogames}{0.41}
\newcommand{\ariAEightAeNogames}{0.38}
\newcommand{\kAAeNogames}{3}
\newcommand{\ariAAeDebut}{0.39}
\newcommand{\ariAEightAeDebut}{0.37}
\newcommand{\kAAeDebut}{3}
\newcommand{\ariAAeConsecutive}{0.62}
\newcommand{\ariAEightAeConsecutive}{0.38}
\newcommand{\kAAeConsecutive}{3}
\newcommand{\ariAAePrimaryBpm}{0.32}
\newcommand{\ariAEightAePrimaryBpm}{0.24}
\newcommand{\kAAePrimaryBpm}{3}
\newcommand{\ariSSummaryOriginal}{0.31}
\newcommand{\ariSEightSummaryOriginal}{0.51}
\newcommand{\kSSummaryOriginal}{3}
\newcommand{\ariSSummaryMp}{0.48}
\newcommand{\ariSEightSummaryMp}{0.59}
\newcommand{\kSSummaryMp}{3}
\newcommand{\ariSSummaryMinSeasons}{0.44}
\newcommand{\ariSEightSummaryMinSeasons}{0.74}
\newcommand{\kSSummaryMinSeasons}{3}
\newcommand{\ariSSummaryNogames}{0.92}
\newcommand{\ariSEightSummaryNogames}{0.92}
\newcommand{\kSSummaryNogames}{4}
\newcommand{\ariSSummaryDebut}{0.41}
\newcommand{\ariSEightSummaryDebut}{0.82}
\newcommand{\kSSummaryDebut}{3}
\newcommand{\ariSSummaryConsecutive}{0.92}
\newcommand{\ariSEightSummaryConsecutive}{0.88}
\newcommand{\kSSummaryConsecutive}{4}
\newcommand{\ariSSummaryPrimaryBpm}{0.29}
\newcommand{\ariSEightSummaryPrimaryBpm}{0.27}
\newcommand{\kSSummaryPrimaryBpm}{3}
\newcommand{\ariSSummaryWard}{0.33}
\newcommand{\ariSEightSummaryWard}{0.53}
\newcommand{\kSSummaryWard}{3}
\newcommand{\ariSMin}{0.29}
\newcommand{\ariSMax}{0.92}
\newcommand{\ariSEightMin}{0.27}
\newcommand{\ariSEightMax}{0.92}
\newcommand{\ariAMin}{0.20}
\newcommand{\ariAMax}{0.62}
\newcommand{\xSummaryRaw}{0.32}
\newcommand{\xSummaryPca}{0.18}
\newcommand{\xSummaryDtw}{0.19}
\newcommand{\xSummaryAe}{0.22}
\newcommand{\xSummaryAeSeedOnly}{0.20}
\newcommand{\xSummaryAeUmap}{0.19}
\newcommand{\xRawPca}{0.12}
\newcommand{\xRawDtw}{0.27}
\newcommand{\xRawAe}{0.20}
\newcommand{\xRawAeSeedOnly}{0.18}
\newcommand{\xRawAeUmap}{0.15}
\newcommand{\xPcaDtw}{0.06}
\newcommand{\xPcaAe}{0.12}
\newcommand{\xPcaAeSeedOnly}{0.12}
\newcommand{\xPcaAeUmap}{0.11}
\newcommand{\xDtwAe}{0.12}
\newcommand{\xDtwAeSeedOnly}{0.13}
\newcommand{\xDtwAeUmap}{0.04}
\newcommand{\xAeAeSeedOnly}{0.57}
\newcommand{\xAeAeUmap}{0.29}
\newcommand{\xAeSeedOnlyAeUmap}{0.20}
\newcommand{\ariAeSummCoarse}{0.25}
\newcommand{\ariAeSummCoarseThree}{0.27}
\newcommand{\nMatchedOld}{1195}
\newcommand{\ariOldEight}{0.17}
\newcommand{\ariOldFour}{0.11}
\newcommand{\ariOldAe}{0.07}
\newcommand{\amiOldEight}{0.29}
\newcommand{\gapSummNormalMax}{0.16}
\newcommand{\gapRawNormalMax}{0.01}
\newcommand{\gapAeNormalMax}{0.11}
\newcommand{\gapSummUniformMax}{1.56}
\newcommand{\dipPmin}{0.71}
\newcommand{\dipPmax}{1.00}
\newcommand{\dipPseasons}{0.000}
\newcommand{\nullSummSilFour}{0.93}
\newcommand{\nullSummSilEight}{0.56}
\newcommand{\nullSummPacEight}{0.22}
\newcommand{\nullSummStableEight}{89}
\newcommand{\copSummSilFour}{0.90}
\newcommand{\copSummSilEight}{0.75}
\newcommand{\copSummStableEight}{96}
\newcommand{\nullRawSilEight}{0.63}
\newcommand{\copRawSilEight}{0.66}
\newcommand{\nullSentence}{A multivariate-normal reference with the same covariance yields a consensus Silhouette of 0.93 at $k=4$, indistinguishable from the real data (0.93), and 0.56 at $k=8$ (real data 0.90). A Gaussian-copula reference, which keeps every real marginal distribution (the integer count of seasons, the bounded shares, the skewed peak values) and the rank correlations but imposes a unimodal dependence structure, yields 0.90 at $k=4$ and 0.75 at $k=8$; for the resampled trajectory the corresponding values are 0.63 (Gaussian) and 0.66 (copula) against 0.85 for the real data. The reproducibility of the four-group partition is therefore fully reproduced by an elongated Gaussian cloud, and most of the reproducibility of the eight-group partition by a unimodal reference that shares only the marginal shapes and rank correlations of the data.}
\newcommand{\nullDiscussion}{structureless references reproduce the consensus stability of the four-group partition entirely and most of that of the eight-group partition.}
\newcommand{\dtwSentence}{reaches a consensus Silhouette of 0.76 at its selected $k=3$ and 0.36 at $k=8$, with $R^2_{\mathrm{shape}}$ 0.66; it is less reproducible than the learned representation and less reproducible than the summary statistics.}
\newcommand{\xMin}{0.04}
\newcommand{\xMax}{0.57}
\newcommand{\eightLotteryA}{71}
\newcommand{\eightStabMin}{0.90}
\newcommand{\eightASmaxOther}{28}
\newcommand{\kSMin}{3}
\newcommand{\kSMax}{4}
\newcommand{\ariSEightBpm}{0.27}
\newcommand{\ariSBpm}{0.29}
\newcommand{\ariSEightMinNoBpm}{0.51}
\newcommand{\ariSEightMaxNoBpm}{0.92}
\newcommand{\ariSEightMinRestrict}{0.51}
\newcommand{\ariSEightMaxRestrict}{0.82}
\newcommand{\corrPerBpm}{0.84}
\newcommand{\ariASeedZero}{0.58}
\newcommand{\ariAdEight}{0.33}
\newcommand{\ariAdOneTwoEight}{0.62}
\newcommand{\ariAdThirtytwo}{0.54}
\newcommand{\ariAdSixteen}{0.29}
\newcommand{\ariABpm}{0.32}
\newcommand{\pcOneVar}{74}
\newcommand{\pcTwoVar}{11}
\newcommand{\pcThreeVar}{6}
\newcommand{\pcOneMean}{1.00}
\newcommand{\pcOneLen}{0.54}
\newcommand{\pcOnePeak}{0.93}
\newcommand{\pcOneWs}{0.73}
\newcommand{\pcTwoSlope}{0.64}
\newcommand{\pcTwoPeakpos}{0.60}
\newcommand{\pcTwoMean}{0.03}
\newcommand{\shapeVarOne}{41}
\newcommand{\shapeVarTwo}{25}
\newcommand{\shapeDipPmin}{0.96}
\newcommand{\dipSlopeP}{1.00}
\newcommand{\peakFirstThird}{32}
\newcommand{\peakLastThird}{23}
\newcommand{\shapeK}{3}
\newcommand{\shapeSil}{0.97}
\newcommand{\shapeSilEight}{0.60}
\newcommand{\shapeNA}{613}
\newcommand{\shapePpA}{0.33}
\newcommand{\shapeSeasA}{8}
\newcommand{\shapeNB}{400}
\newcommand{\shapePpB}{0.50}
\newcommand{\shapeSeasB}{11}
\newcommand{\shapeNC}{663}
\newcommand{\shapePpC}{0.70}
\newcommand{\shapeSeasC}{4}
\newcommand{\shapeLenEps}{0.20}
\newcommand{\nThreeWay}{5}
\newcommand{\rmseTestThreeWay}{0.81}
\newcommand{\rmseTestThreeWaySd}{0.05}
\newcommand{\rmseValThreeWay}{0.74}
\newcommand{\ldaAccThree}{76}
\newcommand{\ldaAccTwo}{56}

\newcommand{\point}[2]{\item[\textbf{#1}] \textit{#2}\par\vspace{2pt}}
\newcommand{\resp}{\par\noindent\textbf{Response.} }
\title{Point-by-point summary of changes relative to the JSET submission\\ \large (for the cover letter of the new submission)}
\author{P. Graovac, \DJ. Had\v{z}i Pavlovi\'c, M. Radoji\v{c}i\'c, S. Radovanovi\'c}
\date{September 2026}
\begin{document}\maketitle
\noindent The manuscript \emph{Discovering NBA Career Trajectories from Longitudinal Performance Data} was declined by \emph{Proc IMechE Part P: Journal of Sports Engineering and Technology} after editorial assessment (29 August 2026) with 26 required concerns and an invitation to resubmit as a new manuscript if a stronger fit could be established. Because the contribution is analytical rather than engineering, the revised study is submitted to a quantitative sports-analytics journal. It is a new study: new data construction, a new validation design and a different substantive conclusion. Below we list every concern and what was done.

\begin{description}[leftmargin=1.2cm,style=nextline]
\point{1}{Clarify the specific contribution.}
\resp The contribution is now explicitly methodological and empirical (Section 1): a reproducibility-oriented protocol for trajectory typing (held-out selection, seed/subsample/hyper-parameter stability, consensus clustering with a prespecified rule, equal-terms baselines, external validation) and its application to complete NBA careers. The journal was changed accordingly.

\point{2}{Strengthen the literature review and define the gap.}
\resp Section 2 now reviews aging-curve models, functional-data trajectory typologies (Wakim and Jin 2014; Elijah et al. 2025), group-based trajectory modelling, sports-economics work on draft and career length, time-series clustering and deep representation learning (including the comparative study of Lafabregue et al. 2022), the documented fragility of UMAP/HDBSCAN pipelines (Chari and Pachter 2023), and consensus/stability criteria (Monti et al. 2003; \c{S}enbabao\u{g}lu et al. 2014; von Luxburg 2010). The gap is stated in the closing paragraph of Section 2.

\point{3}{State explicit research questions.}
\resp Five research questions with prespecified success criteria are stated in Section 1 and answered one by one in Sections 5 and 6.

\point{4}{Complete details of the autoencoder.}
\resp Section 4.1 and Table 1 of the manuscript give architecture, initialisation, optimiser (AdamW), learning rate, weight decay, batch size, gradient clipping, learning-rate schedule, maximum epochs, early-stopping rule and patience, validation split, seeds, loss implementation (Eq.~2) and monitoring; training curves and a 70/10/20 split check are in Appendix A.

\point{5}{Handling of variable-length sequences, gaps, interrupted careers.}
\resp Section 3.3: calendar-axis representation with an observed-indicator, key-padding masks in attention, masked loss, no interpolation; \pctGapPlayers\% of players have at least one gap. Gap deletion is a sensitivity analysis (Section 5.9).

\point{6}{Normalisation of PER.}
\resp Section 3.3: z-scoring with development-set constants, inverse-transformed for reporting; justification given (PER already league-normalised, near-symmetric).

\point{7}{Justify the 64-dimensional latent space.}
\resp Latent dimensionalities 8, 16, 32, 64 and 128 were trained (three to five seeds each). Table 3 reports reconstruction for each; Section 5.9 reports the agreement of the consensus partitions across dimensionalities.

\point{8}{Complete UMAP/HDBSCAN grid.}
\resp Section 4.2 lists all values (\nGrid{} configurations); the full grid results are released as a CSV file.

\point{9}{Justify the composite selection metric.}
\resp The composite is no longer used for selection. Selection uses a prespecified density-based criterion (DBCV with constraints) for the single-run design and a fixed consensus-Silhouette rule for the consensus design (the PAC rule initially specified proved uninformative, which is documented); the old composite is recomputed only to show that different criteria choose different solutions (Appendix Table 10).

\point{10}{Selection bias / overfitting.}
\resp Players were split 80/20 before any modelling. All hyper-parameter selection used the development set; selected configurations were applied to held-out players by UMAP transform and HDBSCAN approximate prediction and evaluated there (Appendix Table 10). The consensus design uses subsampling by construction.

\point{11}{Sensitivity and stability analyses.}
\resp Section 5.2 quantifies single-run reproducibility across UMAP seeds, autoencoder seeds, extraction methods and configurations (ARI, range of $k$, noise, largest-cluster share). Section 5.9 covers latent dimensionality, six sample specifications, gap alignment and the performance metric. This analysis changed the conclusion of the study: the 19-cluster solution is not reproducible.

\point{12}{Baseline comparisons.}
\resp Summary-statistic $k$-means and Ward, PCA $k$-means, DTW $k$-means and raw-trajectory $k$-means are run under the identical consensus protocol and scored on identical criteria (Table 4).

\point{13}{Stronger evaluation of the autoencoder.}
\resp Held-out RMSE is compared with global-mean, player-mean and player-linear-trend baselines under a prespecified criterion (Table 3).

\point{14}{Substantive meaning of clusters, external variables.}
\resp Sections 5.5 and 5.6 test for discrete structure with the gap statistic, the dip test, structureless reference data (Gaussian and copula nulls) and a shape-only analysis, and propose continuous career coordinates as the alternative summary; Section 5.9: position, draft category, All-Star and All-NBA selection, career Win Shares, mean BPM and career games, with Cram\'er's $V$, adjusted mutual information and Kruskal--Wallis $\varepsilon^2$.

\point{15}{Reliance on PER.}
\resp Section 5.8 relates the types to WS and BPM; Section 5.9 repeats the entire pipeline on BPM trajectories and reports the agreement; the limitations of PER are analysed in Sections 6 and 7.

\point{16}{Justify sample restrictions.}
\resp Section 3.2 justifies each restriction; six alternative specifications are evaluated (Table 2, Section 5.9).

\point{17}{Data source, completeness, reproducibility.}
\resp Data are now taken from Basketball-Reference (reference PER implementation) with retrieval date, scraping script, raw download, processed files and full code released (Section 3.1; Data and code availability).

\point{18}{Distinguish results from interpretation.}
\resp Results (Section 5) contain no causal language; Section 6 labels every mechanism as a hypothesis; Section 7 states that no such variable entered the analysis.

\point{19}{Subjective labels and famous players.}
\resp A fixed labelling rule (level term, duration term, slope qualifier with numeric cut-points) is stated before the results; named players are medoid-nearest examples and are explicitly not treated as validation.

\point{20}{Dedicated Limitations section.}
\resp Section 7 addresses every item listed by the Editor.

\point{21}{Narrow the practical claims.}
\resp Claims about scouting, development and roster management are removed; Section 6 states what the typology can and cannot support.

\point{22}{Mathematical notation.}
\resp The manuscript is now typeset in \LaTeX; every equation and symbol was checked in the compiled PDF.

\point{23}{Figures and tables.}
\resp The PCA projection is replaced by figures that show the space actually clustered or the consensus matrix; every figure has labels, legends and uncertainty (interquartile bands, stability scores).

\point{24}{References.}
\resp All references were re-verified (DOIs added; the NFL-draft reference was replaced by NBA-draft studies, Berri et al. 2011 and Coates and Oguntimein 2010); author names with diacritics are handled by Bib\TeX.

\point{25}{Formatting.}
\resp Double-spaced, single-column, 12 pt, in line with the target journal.

\point{26}{Language and AI disclosure.}
\resp The text was rewritten; the AI-assistance statement follows the target journal's policy (Acknowledgements).
\end{description}
\end{document}
